\chapter{Introduction}
\label{ch:introduction}

\section{The aggregation problem}

Several news stories can appear about the same company in one trading session,
but there is only one return observation for that firm and session. Comparing
news with returns therefore requires one daily news measure. Different summaries
retain different information.

The usual summary is the arithmetic mean of the sentiment scores. It measures
the average tone, but opposing stories can cancel out. This dissertation instead
tests
\emph{negative-story share}: the fraction of stories whose strongest FinBERT
label is negative. The share measures how much of the day's coverage is
negative. The mean measures the tone of the average story.

Here, ``hard'' means that each story is assigned to the class with the highest
FinBERT probability. It does not mean that the analysis estimates a confidence
cut-off.

Consider four hard labels: $(-1,-1,+1,+1)$. Their mean is zero, while their
negative-story share is one half. Replacing the two positive labels with neutral
labels leaves the share unchanged but makes the mean more negative. Four mildly
negative stories, by contrast, produce a share of one even if every continuous
score is only slightly below zero. The statistics are related but not
interchangeable.

Earlier research links negative financial language to returns, earnings
information and slow price adjustment
\citep{tetlock2007,tetlock2008words,uhl2014,heston2017}.

Research using public news also reports drift after negative coverage and
reversal after large price moves with no identifiable news story
\citep{chan2003news,savor2012information}. Together, these studies justify
comparing daily news summaries. They do not show that negative-story share is
useful. The measure must add information beyond average sentiment and story
count, precede the return, and survive a later test.

Story count also affects interpretation. The training sample contains 512,153
firm-days, but 48.3\% contain one story and the median is two. On a one-story
day, negative-story share is simply a binary FinBERT label. It describes a
within-day distribution only when several stories are present. Results are
therefore reported for all firm-days and for multi-story days.

The statistical and economic questions remain separate. An association can be
precisely estimated yet too small to cover trading costs, and a risk rule may
reduce losses merely by reducing market exposure. Timing also matters. News
published after a return window opens cannot predict that return. The tests
therefore ask whether the association exists, when it appears, whether it
survives later, and whether it supports a practical rule.

\section{Research questions and hypotheses}

The dissertation first studies the training period, then asks whether the
unchanged relationship survives in a later testing period. The later test uses
choices set in advance. FNSPID, a public financial-news and price dataset, is
used in both periods. Reuters news supplied through LSEG provides a smaller
second sample.

\begin{quote}
\textbf{Training period.} After allowing for average sentiment and news volume,
is a higher negative-story-share rank associated with a lower market-adjusted
return rank from the assigned open to the following open? If so, which part of
that return window contains the relationship?

\textbf{Testing period.} Does the unchanged baseline relationship remain
negative in the fixed 2020--2023 FNSPID news-and-price sample?
\end{quote}

Every coefficient compares firms within the same session. Firms are ranked by
market-adjusted return, negative-story share, average sentiment and story count.
The model asks whether a higher negative-share rank is associated with a lower
return rank after controlling for the other news measures. Daily coefficients
are averaged so that heavily covered sessions do not dominate. Later models
also control for recent returns, volatility and market sensitivity.

A negative coefficient means that firms with more widespread negative coverage
tend to rank lower by return. It does not imply that a given increase in
negative share causes a fixed price fall. Unobserved news or market conditions
may affect both variables, so the result remains a conditional association.

Four hypotheses carry the tests, with the first split into a training-period leg and
a testing-period leg.

\begin{enumerate}[leftmargin=1.1cm]
  \item[H1a.] In the training period, the negative-share coefficient is below zero after controlling for average sentiment and story count.
  \item[H1b.] The unchanged baseline coefficients for all firm-days and multi-story days remain below zero during 2020--2023 after correcting for the two tests.
  \item[H2.] A simple high-minus-low negative-share sort within average-sentiment groups produces a negative training-period spread.
  \item[H3.] A daily long-short rule earns a positive average return after charging 10 basis points per side on its own turnover.
  \item[H4.] After costs, an LSEG risk rule passes its pre-specified return, risk and stability tests against a training-selected control with the same average market exposure.
\end{enumerate}

A high-minus-low sort compares firms with high and low negative-story shares.
Turnover is the amount of the portfolio that must be traded. Ten basis points is
0.10\% of the amount traded on each side.

H1a is exploratory because the nine-rule training comparison selected negative
share. The correction covers the nine initial tests, but the usual interval
around the winning estimate does not capture all selection uncertainty.

H1b is narrower. The two coefficients, the multiple-testing adjustment and the
decision rule were fixed before the 2020--2023 estimates were calculated. The
test is therefore more informative than the training comparison. However,
some portfolio outcomes from those years had already been examined, so the
period was not completely unseen.
\ifdefined\withoutappendices
\else
  \Cref{app:reproducibility} records the origin and run status of every notebook
  used in the chapters that follow.
\fi

\section{Summary of findings}

Only negative-story share remained statistically significant after correcting
the nine training-period comparisons. Firms with a higher negative share tended
to rank lower by assigned-window open-to-next-open abnormal return. After
controlling for average sentiment and story count, the coefficient was
$-0.00831$, with a 95\% confidence interval of $[-0.01410,-0.00252]$ that allows
for dependence between nearby sessions.

The timing checks weaken a forecasting interpretation. The relationship appears
from the assigned open to close, disappears overnight, and is also present from
the previous open to the assigned open. FNSPID lacks reliable row-level
availability times, so some stories may have entered a session after its return
began. The result is a relationship within the assigned window, not evidence of
a forecast based on news known before the return began.

The relationship also fails to persist. In the fixed 2020--2023 testing period,
the unchanged coefficient changes sign to $+0.00583$. A direct comparison finds
a statistically significant difference between the training and testing
estimates. Aggregate data cannot distinguish among changes in timing, coverage,
firm composition or classifier outputs.

The economic evidence also argues against practical usefulness. Across the
observed negative-share rank spread, the full-control estimate corresponds to
$-0.512$ return-rank percentile points. The best break-even cost among the nine
rules is $0.625$ basis points per side, versus the assumed 10 basis points.

The separate Reuters/LSEG samples are too small to show that the relationship
appears in another source. No risk or prompt rule passes its pre-specified
criteria. The training coefficient is therefore a statistical pattern, not a
reliable trading signal.

Chapter 2 reviews the literature, Chapters 3 and 4 describe the data and
methods, Chapter 5 reports the results, and Chapter 6 concludes.
